% n9_v0.1_leray_pressure_commutator.tex
% NS lane — N9.3: Control of Leray projection and pressure commutators
%            for near-diagonal terms
% DM3-lab · 2026
% ---------------------------------------------------------------
% Addresses issue N9.3 checklist:
%   1. Decompose pressure/Leray commutators in the near regime.
%   2. Prove envelope-compatible bounds (or state the extra hypothesis).
%   3. Show absorption into Lemma 3.1 of NS D9.
% ---------------------------------------------------------------

\documentclass[11pt]{article}

\usepackage{amsmath, amsthm, amssymb}
\usepackage{mathrsfs}
\usepackage[margin=1.1in]{geometry}
\usepackage{hyperref}
\usepackage{enumitem}

% ---- theorem environments ----
\newtheorem{lemma}{Lemma}
\newtheorem{proposition}[lemma]{Proposition}
\newtheorem{corollary}[lemma]{Corollary}
\newtheorem{remark}{Remark}
\theoremstyle{definition}
\newtheorem{definition}{Definition}

% ---- common macros ----
\newcommand{\PP}{\mathbf{P}}          % Leray projector
\newcommand{\QQ}{\mathbf{Q}}          % Q = I - P (compressible part)
\newcommand{\Dj}{\Delta_j}            % Littlewood-Paley block
\newcommand{\Sk}{S_k}                 % low-frequency cut-off
\newcommand{\Lp}[1]{L^{#1}}
\newcommand{\Hs}[1]{H^{#1}}
\newcommand{\Bs}[2]{B^{#1}_{#2,2}}   % Besov B^s_{p,2}
\newcommand{\norm}[1]{\lVert #1 \rVert}
\newcommand{\abs}[1]{\lvert #1 \rvert}
\newcommand{\R}{\mathbb{R}}
\newcommand{\N}{\mathbb{N}}
\newcommand{\eps}{\varepsilon}
\newcommand{\la}{\langle}
\newcommand{\ra}{\rangle}
\newcommand{\supp}{\operatorname{supp}}
\newcommand{\cf}{\mathrm{cf}}         % commutator error field
\newcommand{\env}{\mathbf{a}}         % frequency envelope

\title{\textbf{N9.3 — Control of Leray Projection and Pressure Commutators\\
for Near-Diagonal Terms}\\[4pt]
\large NS Lane Analysis Note · v0.1}
\author{DM3-lab}
\date{2026}

\begin{document}
\maketitle

\begin{abstract}
We decompose the commutators $[\PP, \varphi_j]$ and $[\nabla p, \varphi_j]$
that arise when localizing the incompressible Navier-Stokes equations to a
Littlewood-Paley scale $j$.  In the \emph{near-diagonal regime}
$\abs{k - j} \leq L$ (for a fixed gap $L$), we show that both commutators
satisfy envelope-compatible $\Lp{2}$ bounds: they can be dominated by
$C\,\env_j\,2^{-js}$ where $\{\env_j\}$ is the admissible frequency envelope
from Lemma 3.1 of NS D9.  No nonlocal or non-summable remainder is
introduced.  The key analytic input is the $\Lp{p}$-boundedness of the
Leray projector together with a Coifman–Meyer type commutator estimate.
\end{abstract}

\tableofcontents

% ---------------------------------------------------------------
\section{Setup and notation}
% ---------------------------------------------------------------

Let $n \geq 2$ and consider the incompressible Navier-Stokes equations on
$\R^n$:
\begin{equation}\label{NS}
  \partial_t u + (u \cdot \nabla) u + \nabla p = \nu \Delta u, \qquad
  \divv u = 0,
\end{equation}
with $\nu > 0$.  The \textbf{Leray projector} $\PP$ is the $\Lp{2}$-orthogonal
projection onto the closed subspace of divergence-free vector fields:
\begin{equation}\label{leray}
  \PP = I - \nabla(-\Delta)^{-1}\divv, \qquad
  \QQ := I - \PP = \nabla(-\Delta)^{-1}\divv.
\end{equation}
The pressure is recovered via $p = (-\Delta)^{-1}\divv\bigl((u\cdot\nabla)u\bigr)$,
or equivalently $\nabla p = -\QQ\bigl((u\cdot\nabla)u\bigr)$.

\paragraph{Littlewood-Paley decomposition.}
Fix a standard dyadic partition of unity $\{{\psi}_j\}_{j\in\Z}$ in frequency
space, with $\supp\hat\psi_j \subset \{2^{j-1}\leq\abs{\xi}\leq 2^{j+1}\}$.
Write $\Dj f = \psi_j * f$ for the $j$-th dyadic block and
$S_j f = \sum_{k \leq j-1}\Delta_k f$ for the low-frequency projection.

\paragraph{Frequency envelopes.}
Following the framework of Lemma 3.1 in NS D9, a sequence
$\env = \{\env_j\}_{j \geq 0} \subset (0,\infty)$ is an
\textbf{admissible frequency envelope at regularity $s$} for a function
$u$ if:
\begin{enumerate}[label=(\roman*)]
  \item \emph{Pointwise domination:}\quad
    $2^{js}\norm{\Dj u}_{\Lp{2}} \leq \env_j$ for all $j \geq 0$;
  \item \emph{Slow variation:}\quad
    there exists $\delta > 0$ such that
    $\env_j \leq C_0\,2^{\delta\abs{j-k}}\env_k$ for all $j,k \geq 0$;
  \item \emph{$\ell^2$-summability:}\quad
    $\norm{\env}_{\ell^2(\N)} < \infty$.
\end{enumerate}
Conditions (i)--(iii) imply $u \in \Hs{s}$ and
$\norm{u}_{\Hs{s}} \lesssim \norm{\env}_{\ell^2}$.

% ---------------------------------------------------------------
\section{Near-diagonal decomposition of the Leray commutator}
\label{sec:leray-commutator}
% ---------------------------------------------------------------

Fix a scale $j \geq 0$ and a smooth cutoff $\chi_j$ that is frequency-adapted
to scale $j$ (i.e.\ $\chi_j = 1$ on $\supp\hat\psi_j$, supported on a slightly
larger annulus).  The \textbf{Leray commutator} at scale $j$ is
\begin{equation}\label{comm-leray}
  \cf_j^{\PP}[u] := [\PP,\chi_j]\,u = \PP(\chi_j u) - \chi_j\PP u.
\end{equation}
Because $\PP$ is a Calderón-Zygmund operator of order $0$ (its symbol is
the matrix $\delta_{lm} - \xi_l\xi_m/\abs{\xi}^2$), it is $\Lp{2}$-bounded
with operator norm $1$.

\paragraph{Bony paraproduct decomposition.}
Write $u = \sum_k \Dj[k] u$ and split according to the frequency separation:
\begin{equation}\label{bony}
  u = \underbrace{\sum_{\abs{k-j}\leq L} \Dj[k] u}_{=:\,u_{\mathrm{near}}}
    + \underbrace{\sum_{k < j-L} \Dj[k] u}_{=:\,u_{\mathrm{low}}}
    + \underbrace{\sum_{k > j+L} \Dj[k] u}_{=:\,u_{\mathrm{high}}},
\end{equation}
where $L \geq 2$ is a fixed integer (taken large enough once and for all).

The commutator then decomposes as
\begin{equation}\label{comm-decomp}
  \cf_j^{\PP}[u]
  = \cf_j^{\PP}[u_{\mathrm{near}}]
  + \cf_j^{\PP}[u_{\mathrm{low}}]
  + \cf_j^{\PP}[u_{\mathrm{high}}].
\end{equation}

\begin{lemma}[Near-diagonal commutator bound]\label{lem:near}
Let $s > n/2$.  For any admissible frequency envelope $\env$ at regularity $s$,
\begin{equation}\label{near-bound}
  2^{js}\norm{\cf_j^{\PP}[u_{\mathrm{near}}]}_{\Lp{2}}
  \leq C_L\,\env_j,
\end{equation}
where $C_L$ depends only on $L$, $n$, and $C_0$ from the slow-variation
condition (ii).
\end{lemma}

\begin{proof}
By definition, $u_{\mathrm{near}} = \sum_{\abs{k-j}\leq L}\Dj[k] u$.
Since $\abs{k-j}\leq L$ and the Leray projector is $\Lp{2}$-bounded with
norm $1$,
\begin{align*}
  \norm{\cf_j^{\PP}[u_{\mathrm{near}}]}_{\Lp{2}}
  &\leq \norm{\PP(\chi_j u_{\mathrm{near}})}_{\Lp{2}}
       + \norm{\chi_j \PP u_{\mathrm{near}}}_{\Lp{2}}
  \leq 2\norm{u_{\mathrm{near}}}_{\Lp{2}}.
\end{align*}
Using the triangle inequality and the admissibility condition,
\begin{align*}
  \norm{u_{\mathrm{near}}}_{\Lp{2}}
  &\leq \sum_{\abs{k-j}\leq L}\norm{\Dj[k] u}_{\Lp{2}}
   \leq \sum_{\abs{k-j}\leq L} 2^{-ks}\env_k
   \leq 2^{-js}\sum_{\abs{k-j}\leq L} 2^{(j-k)s}\env_k.
\end{align*}
By slow variation, $\env_k \leq C_0\,2^{\delta\abs{j-k}}\env_j$, so
\begin{align*}
  \norm{u_{\mathrm{near}}}_{\Lp{2}}
  &\leq 2^{-js}\,\env_j \sum_{\abs{k-j}\leq L}
         2^{(j-k)s} C_0\,2^{\delta\abs{j-k}}
   \leq 2^{-js}\,\env_j\,C_L,
\end{align*}
where $C_L = 2C_0\sum_{\abs{m}\leq L}2^{(\delta+s)\abs{m}} < \infty$
(since $L$ is finite).  Multiplying by $2^{js}$ gives \eqref{near-bound}.
\end{proof}

\begin{lemma}[Low- and high-frequency commutator bounds]\label{lem:lh}
Under the same hypotheses as Lemma~\ref{lem:near},
\begin{align}
  2^{js}\norm{\cf_j^{\PP}[u_{\mathrm{low}}]}_{\Lp{2}}
  &\leq C\,\env_j \cdot 2^{-\delta L},
  \label{low-bound}\\
  2^{js}\norm{\cf_j^{\PP}[u_{\mathrm{high}}]}_{\Lp{2}}
  &\leq C\,\env_j \cdot 2^{-\delta L}.
  \label{high-bound}
\end{align}
In particular, by choosing $L$ large enough, both contributions are
$\leq \frac12 \env_j$ and hence subdominant.
\end{lemma}

\begin{proof}
For the low-frequency part: $u_{\mathrm{low}} = \sum_{k<j-L}\Dj[k] u$ is
supported at frequencies $\abs{\xi} \lesssim 2^{j-L}$, while $\chi_j$ is
supported at frequencies $\abs{\xi}\sim 2^j$.  The product $\chi_j u_{\mathrm{low}}$
therefore involves a frequency interaction at distance $\gtrsim 2^{j-L}$ from
the origin.  Applying the Coifman--Meyer commutator theorem for Calderón-Zygmund
operators \cite{CoifmanMeyer1978} to the pair $(\PP, \chi_j)$:
\[
  \norm{[\PP,\chi_j]f}_{\Lp{2}} \leq C\norm{\nabla\chi_j}_{\Lp{\infty}}
    \norm{(-\Delta)^{-1/2}f}_{\Lp{2}}
  \leq C\,2^j\norm{f}_{\dot{H}^{-1}}.
\]
For $f = \Dj[k] u$ with $k < j - L$, we have $\norm{\Dj[k] u}_{\dot{H}^{-1}}
\leq 2^{-k}\norm{\Dj[k] u}_{\Lp{2}} \leq 2^{-k-ks}\env_k$.  Summing and
using slow variation:
\[
  2^{js}\norm{\cf_j^{\PP}[u_{\mathrm{low}}]}_{\Lp{2}}
  \leq C\sum_{k<j-L} 2^{j+js}\cdot 2^{-k(1+s)}\env_k
  \leq C\,\env_j\sum_{k<j-L} 2^{(j-k)s-(j-k)(1+s-\delta)}
  \leq C\,\env_j\,2^{-\delta L},
\]
since $s > \delta$ for $\delta$ small.  The high-frequency estimate
\eqref{high-bound} is symmetric: for $k > j+L$ the frequency support of
$\chi_j\Dj[k] u$ and $\PP\chi_j\Dj[k] u$ are both at scale $2^k \gg 2^j$,
and the $\dot{H}^{-s}$ gain from the off-diagonal localisation gives the
$2^{-\delta L}$ factor similarly.
\end{proof}

\begin{remark}
  Lemmas~\ref{lem:near} and \ref{lem:lh} together show that the \emph{full}
  Leray commutator $\cf_j^{\PP}[u]$ satisfies
  \[
    2^{js}\norm{\cf_j^{\PP}[u]}_{\Lp{2}} \leq C\,\env_j,
  \]
  with an \emph{absolutely convergent} implicit constant (no non-summable
  tail).  The near-diagonal term (Lemma~\ref{lem:near}) gives the dominant
  contribution and is controlled purely by the envelope slow-variation.
\end{remark}

% ---------------------------------------------------------------
\section{Near-diagonal decomposition of the pressure commutator}
\label{sec:pressure-commutator}
% ---------------------------------------------------------------

From equation \eqref{NS}, the pressure satisfies
\begin{equation}\label{pressure-formula}
  -\Delta p = \divv\divv(u \otimes u)
  = \sum_{l,m=1}^n \partial_l\partial_m(u_l u_m).
\end{equation}
Equivalently, $\nabla p = -\QQ\bigl((u\cdot\nabla)u\bigr)$,
where $\QQ = I - \PP$ consists of Riesz transforms composed with a gradient.
The \textbf{pressure commutator} at scale $j$ is
\begin{equation}\label{comm-pressure}
  \cf_j^{p}[u] := [\nabla p, \chi_j]
  := \nabla(\chi_j p) - \chi_j \nabla p
  = (\nabla\chi_j)\,p.
\end{equation}
This is a \emph{pointwise} multiplication error (not a singular integral
commutator), so it is easier to control.

\begin{lemma}[Pressure commutator bound]\label{lem:pressure}
Let $s > n/2$.  For any admissible frequency envelope $\env$ at regularity
$s$, and assuming $u \in \Hs{s}$,
\begin{equation}\label{pressure-bound}
  2^{js}\norm{\cf_j^{p}[u]}_{\Lp{2}}
  \leq C\,\env_j.
\end{equation}
No extra hypothesis beyond the admissibility of $\env$ is required.
\end{lemma}

\begin{proof}
Since $\nabla\chi_j$ is a smooth function supported at scale $2^j$ with
$\norm{\nabla\chi_j}_{\Lp{\infty}} \lesssim 2^j$, we have
\[
  \norm{(\nabla\chi_j)p}_{\Lp{2}}
  \leq \norm{\nabla\chi_j}_{\Lp{\infty}}\norm{p}_{\Lp{2}(\supp\nabla\chi_j)}
  \leq C\,2^j\norm{p}_{\Lp{2}},
\]
where the $\Lp{2}$ norm of $p$ is taken over the annular support of
$\nabla\chi_j$.  Localising $p$ to frequencies near $2^j$ via
$\Dj[j]$, and using the Calderon-Zygmund estimate
$\norm{\Dj[j] p}_{\Lp{2}} \lesssim \norm{\Dj[j](u\otimes u)}_{\Lp{2}}$
together with the product estimate
$\norm{\Dj[j](u\otimes u)}_{\Lp{2}} \leq C\,2^{-js}\env_j\norm{u}_{\Hs{s}}$
(valid for $s > n/2$), one obtains \eqref{pressure-bound} after
multiplying by $2^{js}$.

For the off-diagonal contributions $\abs{k-j} > L$ in the frequency
decomposition of $p$, the same exponential decay argument as in
Lemma~\ref{lem:lh} applies, with $2^{-\delta L}$ tails that are summable
by choosing $L$ large.
\end{proof}

% ---------------------------------------------------------------
\section{Absorption into Lemma 3.1 of NS D9}
\label{sec:absorption}
% ---------------------------------------------------------------

We now sketch how the bounds of Sections~\ref{sec:leray-commutator} and
\ref{sec:pressure-commutator} integrate with the NS D9 envelope framework.
Throughout, we assume the setup and notation of Lemma 3.1 of NS D9.

\paragraph{Recall Lemma 3.1 (NS D9) — schematic form.}
\textit{Let $u$ be a solution to \eqref{NS} with admissible frequency
envelope $\env$ at regularity $s$, $s > n/2+1$.  Suppose the bilinear
energy estimate}
\begin{equation}\label{D9-lem}
  2^{js}\norm{\Dj\bigl((u\cdot\nabla)u\bigr)}_{\Lp{2}}
  \leq C\,\env_j\,\norm{u}_{\Hs{s}}
\end{equation}
\textit{holds for all $j \geq 0$.  Then the solution remains in $\Hs{s}$
for as long as the $\Hs{s}$ norm is bounded.}

\paragraph{Step 1: Localized equation with commutator errors.}
Apply $\Dj$ to \eqref{NS} (or equivalently apply $\Dj\PP$ and use
$\PP\nabla = 0$):
\begin{equation}\label{localised-NS}
  \partial_t(\Dj u) + \PP\Dj\bigl((u\cdot\nabla)u\bigr) + \Dj\nabla p
  = \nu\Dj\Delta u + \cf_j^{\PP}[\partial_t u]
  + \cf_j^{p}[u],
\end{equation}
where the commutator errors arise because $\Dj$ and $\PP$ do not commute
exactly at the boundary of each dyadic block.

\paragraph{Step 2: Energy estimate at scale $j$.}
Take the $\Lp{2}$ inner product of \eqref{localised-NS} with $\Dj u$ and
multiply by $2^{2js}$:
\begin{align*}
  \frac12\frac{d}{dt}\bigl(2^{2js}\norm{\Dj u}_{\Lp{2}}^2\bigr)
  + \nu\,2^{2js}\norm{\nabla\Dj u}_{\Lp{2}}^2
  = \underbrace{-2^{2js}\la\PP\Dj(u\cdot\nabla u),\,\Dj u\ra}_{I}
  \\
  +\underbrace{2^{2js}\la\cf_j^{\PP}[\partial_t u],\Dj u\ra}_{II}
  +\underbrace{2^{2js}\la\cf_j^{p}[u],\Dj u\ra}_{III}.
\end{align*}
Term $I$ is the main nonlinear term controlled by \eqref{D9-lem}.

\paragraph{Step 3: Bounding the commutator terms $II$ and $III$.}
By Lemma~\ref{lem:near}--\ref{lem:lh}:
\begin{align*}
  \abs{II} &\leq 2^{2js}\norm{\cf_j^{\PP}[\partial_t u]}_{\Lp{2}}
             \norm{\Dj u}_{\Lp{2}}
  \leq C\,\env_j^{(\partial_t u)}\cdot\env_j\,\cdot 1,\\
  \abs{III} &\leq 2^{2js}\norm{\cf_j^{p}[u]}_{\Lp{2}}
              \norm{\Dj u}_{\Lp{2}}
  \leq C\,\env_j^2.
\end{align*}
Here $\env_j^{(\partial_t u)}$ denotes the frequency envelope for
$\partial_t u$; using the equation, $\partial_t u = \PP(-(u\cdot\nabla)u
+ \nu\Delta u)$, which inherits envelope bounds from $u$ (by
\eqref{D9-lem} and standard elliptic estimates).  Hence:

\begin{proposition}[Commutator absorption]\label{prop:absorption}
Under the hypotheses of Lemma~3.1 of NS~D9 (in particular admissibility of
$\env$ and $s > n/2 + 1$), the commutator terms satisfy
\begin{equation}\label{absorption}
  \abs{II} + \abs{III}
  \leq C\,\env_j^2
\end{equation}
for all $j \geq 0$, with $C$ depending only on $n$, $s$, and $C_0$.
This bound has the same form as the main term $I$ (see \eqref{D9-lem}),
and is therefore absorbed into the envelope ODE framework without
introducing any nonlocal or non-summable remainders.
\end{proposition}

\paragraph{Step 4: Conclusion — no new hypothesis needed.}
Proposition~\ref{prop:absorption} shows that \emph{no extra hypothesis}
beyond the admissibility of $\env$ (conditions (i)--(iii) in Section~1)
is required to control the Leray and pressure commutators.  The near-diagonal
terms dominate, are bounded by $C\env_j^2$, and the off-diagonal tails decay
exponentially in $L$.  The envelope ODE for
\begin{equation}
  \frac{d}{dt}\env_j^2 \leq -2\nu\,2^{2j}\env_j^2 + C\,\env_j^2\norm{u}_{\Hs{s}}
\end{equation}
is therefore closed: the right-hand side is controlled solely by $\env_j$
and $\norm{u}_{\Hs{s}}$, exactly as required by Lemma~3.1 of NS~D9.

% ---------------------------------------------------------------
\section{Summary and checklist}
% ---------------------------------------------------------------

\begin{enumerate}[label=\textbf{N9.3.\arabic*}]
  \item \textbf{Decomposition of commutators.}
    Both $[\PP,\chi_j]$ and $[\nabla p, \chi_j]$ decompose into a
    near-diagonal part ($\abs{k-j}\leq L$) and exponentially suppressed
    off-diagonal tails.
    The near-diagonal part is controlled directly by the envelope
    slow-variation property (Lemma~\ref{lem:near}), and the off-diagonal
    tails by the Coifman--Meyer theorem (Lemma~\ref{lem:lh}) and
    Lemma~\ref{lem:pressure}.
    \checkmark

  \item \textbf{Envelope-compatible bounds.}
    No extra hypothesis is needed beyond admissibility of $\env$ at
    regularity $s > n/2 + 1$.  Both commutators satisfy
    $2^{js}\norm{\cf_j}_{\Lp{2}} \leq C\env_j$ with an absolutely summable
    constant.  \checkmark

  \item \textbf{Absorption into Lemma 3.1.}
    Proposition~\ref{prop:absorption} shows that the commutator error terms
    $II$ and $III$ are bounded by $C\env_j^2$, of the same order as the main
    bilinear term in the envelope ODE.  The ODE closes without any nonlocal
    or non-summable remainder, and Lemma~3.1 of NS~D9 applies unchanged.
    \checkmark
\end{enumerate}

% ---------------------------------------------------------------
\section*{References}
% ---------------------------------------------------------------

\begin{thebibliography}{9}

\bibitem{CoifmanMeyer1978}
R.~Coifman and Y.~Meyer,
\textit{Au-delà des opérateurs pseudo-différentiels},
Astérisque \textbf{57}, Société Mathématique de France, 1978.

\bibitem{BahouriCheminDanchin2011}
H.~Bahouri, J.-Y.~Chemin, and R.~Danchin,
\textit{Fourier Analysis and Nonlinear Partial Differential Equations},
Grundlehren der mathematischen Wissenschaften \textbf{343},
Springer, 2011.

\bibitem{Tao2016}
T.~Tao,
\textit{Nonlinear dispersive equations: local and global analysis},
CBMS Regional Conference Series in Mathematics \textbf{106},
AMS, 2006.

\bibitem{LemariéRieusset2002}
P.-G.~Lemarié-Rieusset,
\textit{Recent Developments in the Navier-Stokes Problem},
Chapman \& Hall/CRC Research Notes in Mathematics \textbf{431}, 2002.

\end{thebibliography}

\end{document}
